% Options for packages loaded elsewhere
\PassOptionsToPackage{unicode}{hyperref}
\PassOptionsToPackage{hyphens}{url}
\documentclass[
]{article}
\usepackage{xcolor}
\usepackage[margin=1in]{geometry}
\usepackage{amsmath,amssymb}
\setcounter{secnumdepth}{-\maxdimen} % remove section numbering
\usepackage{iftex}
\ifPDFTeX
  \usepackage[T1]{fontenc}
  \usepackage[utf8]{inputenc}
  \usepackage{textcomp} % provide euro and other symbols
\else % if luatex or xetex
  \usepackage{unicode-math} % this also loads fontspec
  \defaultfontfeatures{Scale=MatchLowercase}
  \defaultfontfeatures[\rmfamily]{Ligatures=TeX,Scale=1}
\fi
\usepackage{lmodern}
\ifPDFTeX\else
  % xetex/luatex font selection
\fi
% Use upquote if available, for straight quotes in verbatim environments
\IfFileExists{upquote.sty}{\usepackage{upquote}}{}
\IfFileExists{microtype.sty}{% use microtype if available
  \usepackage[]{microtype}
  \UseMicrotypeSet[protrusion]{basicmath} % disable protrusion for tt fonts
}{}
\makeatletter
\@ifundefined{KOMAClassName}{% if non-KOMA class
  \IfFileExists{parskip.sty}{%
    \usepackage{parskip}
  }{% else
    \setlength{\parindent}{0pt}
    \setlength{\parskip}{6pt plus 2pt minus 1pt}}
}{% if KOMA class
  \KOMAoptions{parskip=half}}
\makeatother
\usepackage{color}
\usepackage{fancyvrb}
\newcommand{\VerbBar}{|}
\newcommand{\VERB}{\Verb[commandchars=\\\{\}]}
\DefineVerbatimEnvironment{Highlighting}{Verbatim}{commandchars=\\\{\}}
% Add ',fontsize=\small' for more characters per line
\usepackage{framed}
\definecolor{shadecolor}{RGB}{248,248,248}
\newenvironment{Shaded}{\begin{snugshade}}{\end{snugshade}}
\newcommand{\AlertTok}[1]{\textcolor[rgb]{0.94,0.16,0.16}{#1}}
\newcommand{\AnnotationTok}[1]{\textcolor[rgb]{0.56,0.35,0.01}{\textbf{\textit{#1}}}}
\newcommand{\AttributeTok}[1]{\textcolor[rgb]{0.13,0.29,0.53}{#1}}
\newcommand{\BaseNTok}[1]{\textcolor[rgb]{0.00,0.00,0.81}{#1}}
\newcommand{\BuiltInTok}[1]{#1}
\newcommand{\CharTok}[1]{\textcolor[rgb]{0.31,0.60,0.02}{#1}}
\newcommand{\CommentTok}[1]{\textcolor[rgb]{0.56,0.35,0.01}{\textit{#1}}}
\newcommand{\CommentVarTok}[1]{\textcolor[rgb]{0.56,0.35,0.01}{\textbf{\textit{#1}}}}
\newcommand{\ConstantTok}[1]{\textcolor[rgb]{0.56,0.35,0.01}{#1}}
\newcommand{\ControlFlowTok}[1]{\textcolor[rgb]{0.13,0.29,0.53}{\textbf{#1}}}
\newcommand{\DataTypeTok}[1]{\textcolor[rgb]{0.13,0.29,0.53}{#1}}
\newcommand{\DecValTok}[1]{\textcolor[rgb]{0.00,0.00,0.81}{#1}}
\newcommand{\DocumentationTok}[1]{\textcolor[rgb]{0.56,0.35,0.01}{\textbf{\textit{#1}}}}
\newcommand{\ErrorTok}[1]{\textcolor[rgb]{0.64,0.00,0.00}{\textbf{#1}}}
\newcommand{\ExtensionTok}[1]{#1}
\newcommand{\FloatTok}[1]{\textcolor[rgb]{0.00,0.00,0.81}{#1}}
\newcommand{\FunctionTok}[1]{\textcolor[rgb]{0.13,0.29,0.53}{\textbf{#1}}}
\newcommand{\ImportTok}[1]{#1}
\newcommand{\InformationTok}[1]{\textcolor[rgb]{0.56,0.35,0.01}{\textbf{\textit{#1}}}}
\newcommand{\KeywordTok}[1]{\textcolor[rgb]{0.13,0.29,0.53}{\textbf{#1}}}
\newcommand{\NormalTok}[1]{#1}
\newcommand{\OperatorTok}[1]{\textcolor[rgb]{0.81,0.36,0.00}{\textbf{#1}}}
\newcommand{\OtherTok}[1]{\textcolor[rgb]{0.56,0.35,0.01}{#1}}
\newcommand{\PreprocessorTok}[1]{\textcolor[rgb]{0.56,0.35,0.01}{\textit{#1}}}
\newcommand{\RegionMarkerTok}[1]{#1}
\newcommand{\SpecialCharTok}[1]{\textcolor[rgb]{0.81,0.36,0.00}{\textbf{#1}}}
\newcommand{\SpecialStringTok}[1]{\textcolor[rgb]{0.31,0.60,0.02}{#1}}
\newcommand{\StringTok}[1]{\textcolor[rgb]{0.31,0.60,0.02}{#1}}
\newcommand{\VariableTok}[1]{\textcolor[rgb]{0.00,0.00,0.00}{#1}}
\newcommand{\VerbatimStringTok}[1]{\textcolor[rgb]{0.31,0.60,0.02}{#1}}
\newcommand{\WarningTok}[1]{\textcolor[rgb]{0.56,0.35,0.01}{\textbf{\textit{#1}}}}
\usepackage{graphicx}
\makeatletter
\newsavebox\pandoc@box
\newcommand*\pandocbounded[1]{% scales image to fit in text height/width
  \sbox\pandoc@box{#1}%
  \Gscale@div\@tempa{\textheight}{\dimexpr\ht\pandoc@box+\dp\pandoc@box\relax}%
  \Gscale@div\@tempb{\linewidth}{\wd\pandoc@box}%
  \ifdim\@tempb\p@<\@tempa\p@\let\@tempa\@tempb\fi% select the smaller of both
  \ifdim\@tempa\p@<\p@\scalebox{\@tempa}{\usebox\pandoc@box}%
  \else\usebox{\pandoc@box}%
  \fi%
}
% Set default figure placement to htbp
\def\fps@figure{htbp}
\makeatother
\setlength{\emergencystretch}{3em} % prevent overfull lines
\providecommand{\tightlist}{%
  \setlength{\itemsep}{0pt}\setlength{\parskip}{0pt}}
\usepackage{bookmark}
\IfFileExists{xurl.sty}{\usepackage{xurl}}{} % add URL line breaks if available
\urlstyle{same}
% fallback for those not using the hyperref driver hyperxmp:
\makeatletter
\@ifundefined{xmpquote}{\newcommand{\xmpquote}[1]{#1}}{}
\makeatother
\hypersetup{
  pdftitle={Diamonds Dataset - Descriptive Statistics},
  pdfauthor={Christine Donovan},
  hidelinks,
  pdfcreator={LaTeX via pandoc}}

\title{Diamonds Dataset - Descriptive Statistics}
\author{Christine Donovan}
\date{2026-09-12}

\begin{document}
\maketitle

\section{This is a heading}\label{this-is-a-heading}

This is regular text/narrative.

\subsection{Code Chunk}\label{code-chunk}

\begin{Shaded}
\begin{Highlighting}[]
\CommentTok{\# R code goes here}
\NormalTok{diamonds }\OtherTok{\textless{}{-}} \FunctionTok{read.csv}\NormalTok{(}\StringTok{"https://raw.githubusercontent.com/chrdonovan{-}hub/64060\_cdonova2/main/Assignment\_1/diamonds.csv"}\NormalTok{)}
\FunctionTok{head}\NormalTok{(diamonds)}
\end{Highlighting}
\end{Shaded}

\begin{verbatim}
##   carat       cut color clarity depth table price    x    y    z
## 1  0.23     Ideal     E     SI2  61.5    55   326 3.95 3.98 2.43
## 2  0.21   Premium     E     SI1  59.8    61   326 3.89 3.84 2.31
## 3  0.23      Good     E     VS1  56.9    65   327 4.05 4.07 2.31
## 4  0.29   Premium     I     VS2  62.4    58   334 4.20 4.23 2.63
## 5  0.31      Good     J     SI2  63.3    58   335 4.34 4.35 2.75
## 6  0.24 Very Good     J    VVS2  62.8    57   336 3.94 3.96 2.48
\end{verbatim}

\begin{Shaded}
\begin{Highlighting}[]
\FunctionTok{summary}\NormalTok{(diamonds)}
\end{Highlighting}
\end{Shaded}

\begin{verbatim}
##      carat               cut            color          clarity   
##  Min.   :0.2000   Length   :101   Length   :101   Length   :101  
##  1st Qu.:0.2400   N.unique :  5   N.unique :  5   N.unique :  6  
##  Median :0.2600   N.blank  :  0   N.blank  :  0   N.blank  :  0  
##  Mean   :0.2636   Min.nchar:  4   Min.nchar:  1   Min.nchar:  3  
##  3rd Qu.:0.3000   Max.nchar:  9   Max.nchar:  1   Max.nchar:  4  
##  Max.   :0.3200                                                  
##      depth          table           price             x               y        
##  Min.   :56.9   Min.   :54.00   Min.   :326.0   Min.   :3.790   Min.   :3.750  
##  1st Qu.:62.8   1st Qu.:56.00   1st Qu.:354.0   1st Qu.:3.970   1st Qu.:4.010  
##  Median :63.1   Median :56.00   Median :360.0   Median :4.080   Median :4.100  
##  Mean   :62.9   Mean   :56.94   Mean   :355.9   Mean   :4.086   Mean   :4.108  
##  3rd Qu.:63.4   3rd Qu.:57.00   3rd Qu.:362.0   3rd Qu.:4.200   3rd Qu.:4.230  
##  Max.   :65.1   Max.   :69.00   Max.   :363.0   Max.   :4.380   Max.   :4.460  
##        z        
##  Min.   :2.270  
##  1st Qu.:2.520  
##  Median :2.590  
##  Mean   :2.576  
##  3rd Qu.:2.670  
##  Max.   :2.810
\end{verbatim}

\begin{Shaded}
\begin{Highlighting}[]
\FunctionTok{library}\NormalTok{(psych)}
\FunctionTok{describe}\NormalTok{(diamonds[, }\FunctionTok{c}\NormalTok{(}\StringTok{"price"}\NormalTok{, }\StringTok{"carat"}\NormalTok{, }\StringTok{"depth"}\NormalTok{, }\StringTok{"table"}\NormalTok{, }\StringTok{"x"}\NormalTok{, }\StringTok{"y"}\NormalTok{, }\StringTok{"z"}\NormalTok{)])}
\end{Highlighting}
\end{Shaded}

\begin{verbatim}
##       vars   n   mean   sd median trimmed  mad    min    max range  skew
## price    1 101 355.89 9.37 360.00  357.81 4.45 326.00 363.00 37.00 -1.63
## carat    2 101   0.26 0.03   0.26    0.26 0.04   0.20   0.32  0.12  0.05
## depth    3 101  62.90 1.15  63.10   63.05 0.44  56.90  65.10  8.20 -2.08
## table    4 101  56.94 2.26  56.00   56.49 0.00  54.00  69.00 15.00  2.55
## x        5 101   4.09 0.14   4.08    4.08 0.16   3.79   4.38  0.59  0.27
## y        6 101   4.11 0.16   4.10    4.11 0.15   3.75   4.46  0.71 -0.07
## z        7 101   2.58 0.12   2.59    2.58 0.12   2.27   2.81  0.54 -0.46
##       kurtosis   se
## price     1.76 0.93
## carat    -1.13 0.00
## depth     7.32 0.11
## table     8.06 0.23
## x        -0.69 0.01
## y        -0.41 0.02
## z        -0.22 0.01
\end{verbatim}

\begin{Shaded}
\begin{Highlighting}[]
\CommentTok{\# Cut quality}
\FunctionTok{table}\NormalTok{(diamonds}\SpecialCharTok{$}\NormalTok{cut)}
\end{Highlighting}
\end{Shaded}

\begin{verbatim}
## 
##      Fair      Good     Ideal   Premium Very Good 
##        16        45         4        15        21
\end{verbatim}

\begin{Shaded}
\begin{Highlighting}[]
\CommentTok{\# Color}
\FunctionTok{table}\NormalTok{(diamonds}\SpecialCharTok{$}\NormalTok{color)}
\end{Highlighting}
\end{Shaded}

\begin{verbatim}
## 
##  E  F  H  I  J 
## 10  3 24 57  7
\end{verbatim}

\begin{Shaded}
\begin{Highlighting}[]
\CommentTok{\# Clarity}
\FunctionTok{table}\NormalTok{(diamonds}\SpecialCharTok{$}\NormalTok{clarity)}
\end{Highlighting}
\end{Shaded}

\begin{verbatim}
## 
##  SI1  SI2  VS1  VS2 VVS1 VVS2 
##   75   15    4    5    1    1
\end{verbatim}

\begin{Shaded}
\begin{Highlighting}[]
\CommentTok{\# Install if not already installed}
\ControlFlowTok{if}\NormalTok{ (}\SpecialCharTok{!}\FunctionTok{require}\NormalTok{(}\StringTok{"psych"}\NormalTok{, }\AttributeTok{quietly =} \ConstantTok{TRUE}\NormalTok{)) \{}
  \FunctionTok{install.packages}\NormalTok{(}\StringTok{"psych"}\NormalTok{)}
\NormalTok{\}}
\FunctionTok{library}\NormalTok{(psych)}

\FunctionTok{library}\NormalTok{(dplyr)}
\end{Highlighting}
\end{Shaded}

\begin{verbatim}
## 
## Attaching package: 'dplyr'
\end{verbatim}

\begin{verbatim}
## The following objects are masked from 'package:stats':
## 
##     filter, lag
\end{verbatim}

\begin{verbatim}
## The following objects are masked from 'package:base':
## 
##     intersect, setdiff, setequal, union
\end{verbatim}

\begin{Shaded}
\begin{Highlighting}[]
\NormalTok{diamonds }\SpecialCharTok{\%\textgreater{}\%}
  \FunctionTok{group\_by}\NormalTok{(cut) }\SpecialCharTok{\%\textgreater{}\%}
  \FunctionTok{summarize}\NormalTok{(}
    \AttributeTok{Count =} \FunctionTok{n}\NormalTok{(),}
    \AttributeTok{Mean\_Price =} \FunctionTok{mean}\NormalTok{(price),}
    \AttributeTok{Median\_Price =} \FunctionTok{median}\NormalTok{(price),}
    \AttributeTok{Std\_Dev =} \FunctionTok{sd}\NormalTok{(price),}
    \AttributeTok{Min\_Price =} \FunctionTok{min}\NormalTok{(price),}
    \AttributeTok{Max\_Price =} \FunctionTok{max}\NormalTok{(price)}
\NormalTok{  )}
\end{Highlighting}
\end{Shaded}

\begin{verbatim}
## # A tibble: 5 x 7
##   cut       Count Mean_Price Median_Price Std_Dev Min_Price Max_Price
##   <chr>     <int>      <dbl>        <dbl>   <dbl>     <int>     <int>
## 1 Fair         16       358.          359    6.16       337       363
## 2 Good         45       358.          361    7.70       327       363
## 3 Ideal         4       340.          342    9.57       326       348
## 4 Premium      15       353.          357   11.5        326       363
## 5 Very Good    21       355.          360    9.98       336       363
\end{verbatim}

\pandocbounded{\includegraphics[keepaspectratio]{}} \# Box plot of price
by cut boxplot(diamonds\$price \textasciitilde{} diamonds\$cut, main =
``Diamond Price by Cut Quality'', xlab = ``Cut'', ylab = ``Price'')

\subsection{Step 4: Run/Test Code
Chunks}\label{step-4-runtest-code-chunks}

\subsection{Step 5: ``Knit'' Your
Document}\label{step-5-knit-your-document}

\subsection{Step 6: Customize Output
(Optional)}\label{step-6-customize-output-optional}

\section{This code runs but output is
hidden}\label{this-code-runs-but-output-is-hidden}

\section{echo=FALSE hides the code}\label{echofalse-hides-the-code}

\section{warning=FALSE and message=FALSE hide
warnings/messages}\label{warningfalse-and-messagefalse-hide-warningsmessages}

summary(diamonds)

Or to show code but not output:

\begin{Shaded}
\begin{Highlighting}[]
\CommentTok{\# This code is shown but NOT executed}
\FunctionTok{install.packages}\NormalTok{(}\StringTok{"tidyverse"}\NormalTok{)}
\end{Highlighting}
\end{Shaded}

\subsection{Pro Tips:}\label{pro-tips}

\begin{itemize}
\tightlist
\item
  \textbf{Save as you go}: Save your \texttt{.Rmd} file frequently
  (Ctrl+S)
\item
  \textbf{Use headings}: \texttt{\#\ Title}, \texttt{\#\#\ Subtitle},
  \texttt{\#\#\#\ Sub-subtitle} organize your report
\item
  \textbf{Inline code}: Write \texttt{355.8910891} in text to embed R
  calculations
\item
  \textbf{Output formats}: Change \texttt{output:\ html\_document} to
  \texttt{output:\ pdf\_document} or \texttt{output:\ word\_document} in
  the header \#\# R Markdown
\end{itemize}

This is an R Markdown document. Markdown is a simple formatting syntax
for authoring HTML, PDF, and MS Word documents. For more details on
using R Markdown see \url{http://rmarkdown.rstudio.com}.

When you click the \textbf{Knit} button a document will be generated
that includes both content as well as the output of any embedded R code
chunks within the document. You can embed an R code chunk like this:

\subsection{Including Plots}\label{including-plots}

You can also embed plots, for example:

\pandocbounded{\includegraphics[keepaspectratio]{Assignment_1_files/figure-latex/pressure-1.pdf}}

Note that the \texttt{echo\ =\ FALSE} parameter was added to the code
chunk to prevent printing of the R code that generated the plot.

\end{document}
